\begin{examplebox}
	\textbf{\textcolor{CExample}{例1（基础）}}

	\textbf{\textcolor{CExample}{题目：}}
	斜线$l$与平面$\alpha$所成角为$30^{\circ}$，$l$在$\alpha$内的射影$l'$与$\alpha$内一直线$m$（过斜足）所成角为$45^{\circ}$，求$l$与$m$所成角$\theta$的余弦值。

	\textbf{\textcolor{CExample}{解题步骤：}}
	\begin{description}[style=nextline, leftmargin=2.8em, labelsep=0.5em, itemsep=0.45em]
		\item[\textbf{第一步（识别角大小）：}] 由题意，$\theta_1 = 30^{\circ}$，$\theta_2 = 45^{\circ}$。
			\par\textbf{理由：} $\theta_1$是线面角，$\theta_2$是射影角。
		\item[\textbf{第二步（套用定理）：}] 根据三余弦定理，$\cos\theta = \cos\theta_1 \cdot \cos\theta_2$。
			\par\textbf{理由：} 最小角定理直接应用。
		\item[\textbf{第三步（计算得结论）：}] \par
			\[
				\begin{aligned}
					\cos\theta &= \cos30^{\circ} \cdot \cos45^{\circ} \\
					&= \frac{\sqrt{3}}{2} \cdot \frac{\sqrt{2}}{2} \\
					&= \frac{\sqrt{6}}{4}.
			\end{aligned}
			\]
			\par\textbf{理由：} 代数计算。
	\end{description}

	\textbf{\textcolor{CExample}{关键结论：}}
	$\cos\theta = \frac{\sqrt{6}}{4}$.

	\medskip
	\textbf{\textcolor{CExample}{例2（稍变形）}}

	\textbf{\textcolor{CExample}{题目：}}
	斜线$l$与平面$\alpha$所成角为$30^{\circ}$，$l$与$\alpha$内一直线$m$（过斜足）所成角为$60^{\circ}$，求射影$l'$与$m$所成角$\theta_2$的余弦值。

	\textbf{\textcolor{CExample}{解题步骤：}}
	\begin{description}[style=nextline, leftmargin=2.8em, labelsep=0.5em, itemsep=0.45em]
		\item[\textbf{第一步（识别已知角）：}] 由题意，$\theta_1 = 30^{\circ}$，$\theta = 60^{\circ}$。
			\par\textbf{理由：} $\theta_1$为线面角，$\theta$为斜线与$m$的夹角。
		\item[\textbf{第二步（变形公式）：}] 由$\cos\theta = \cos\theta_1\cos\theta_2$，得$\cos\theta_2 = \dfrac{\cos\theta}{\cos\theta_1}$。
			\par\textbf{理由：} 等式的变形。
		\item[\textbf{第三步（代入求值）：}] \par
			\[
				\begin{aligned}
					\cos\theta_2 &= \dfrac{\cos60^{\circ}}{\cos30^{\circ}} \\
					&= \dfrac{\frac{1}{2}}{\frac{\sqrt{3}}{2}} \\
					&= \dfrac{1}{\sqrt{3}} \\
					&= \dfrac{\sqrt{3}}{3}.
			\end{aligned}
			\]
			\par\textbf{理由：} 代入特殊角计算，并有理化。
	\end{description}

	\textbf{\textcolor{CExample}{关键结论：}}
	$\cos\theta_2 = \frac{\sqrt{3}}{3}$.

	\medskip
	\textbf{\textcolor{CExample}{例3（综合应用）}}

	\textbf{\textcolor{CExample}{题目：}}
	在棱长为$1$的正方体$ABCD-A_1B_1C_1D_1$中，$BD_1$为体对角线，$AB$是底面内过$B$的棱，求$BD_1$与$AB$所成角（锐角）的余弦值。

	\textbf{\textcolor{CExample}{解题步骤：}}
	\begin{description}[style=nextline, leftmargin=2.8em, labelsep=0.5em, itemsep=0.45em]
		\item[\textbf{第一步（确定各角）：}] 取$D_1$在底面的投影为$D$，则$BD$为$BD_1$在底面上的射影，故$\theta_1 = \angle D_1BD$（线面角），$\theta_2 = \angle ABD$（射影$BD$与$AB$的夹角）。
			\par\textbf{理由：} 由射影定义及线面角定义。
		\item[\textbf{第二步（计算$\cos\theta_1$与$\cos\theta_2$）：}] \par
			在 $\mathrm{Rt}\triangle D_1BD$ 中，$BD = \sqrt{2}$，$BD_1 = \sqrt{3}$，故
			\[
				\begin{aligned}
					\cos\theta_1 &= \dfrac{BD}{BD_1} \\
					&= \dfrac{\sqrt{2}}{\sqrt{3}} \\
					&= \dfrac{\sqrt{6}}{3}.
			\end{aligned}
			\]
			在 $\mathrm{Rt}\triangle ABD$ 中，$\angle ABD = \theta_2$，$AB=1$，$BD=\sqrt{2}$，故
			\[
				\begin{aligned}
					\cos\theta_2 &= \dfrac{AB}{BD} \\
					&= \dfrac{1}{\sqrt{2}} \\
					&= \dfrac{\sqrt{2}}{2}.
			\end{aligned}
			\]
			\par\textbf{理由：} 直角三角形的边角关系。
		\item[\textbf{第三步（套用三余弦定理求$\theta$）：}] \par
			\[
				\begin{aligned}
					\cos\theta &= \cos\theta_1\cos\theta_2 \\
					&= \dfrac{\sqrt{6}}{3} \times \dfrac{\sqrt{2}}{2} \\
					&= \dfrac{\sqrt{12}}{6} \\
					&= \dfrac{2\sqrt{3}}{6} \\
					&= \dfrac{\sqrt{3}}{3}.
			\end{aligned}
			\]
			\par\textbf{理由：} 直接代入三余弦公式。
	\end{description}

	\textbf{\textcolor{CExample}{关键结论：}}
	$BD_1$与$AB$所成角的余弦值为$\dfrac{\sqrt{3}}{3}$.
\end{examplebox}
